\documentclass[11pt]{article}

\usepackage[margin=1in]{geometry}
\usepackage{array}
\usepackage{booktabs}
\usepackage{longtable}
\usepackage{tabularx}
\usepackage{xcolor}
\usepackage{hyperref}

\hypersetup{
  colorlinks=true,
  linkcolor=black,
  urlcolor=blue
}

\setlength{\parindent}{0pt}
\setlength{\parskip}{6pt}
\renewcommand{\arraystretch}{1.18}

\newcolumntype{Y}{>{\raggedright\arraybackslash}X}
\newcommand{\markcell}[2]{\textbf{#1/#2}}

\begin{document}

\begin{center}
  {\Large Geography Paper 2 TZ1 HLSL 2025: Student Work Marking Report}\\[4pt]
  {\normalsize Marked against \texttt{Geography\_paper\_2\_TZ1\_HLSL\_markscheme-2025.pdf}}
\end{center}

\section*{Overall Result}

\begin{tabularx}{\textwidth}{@{}p{0.34\textwidth} c c Y@{}}
\toprule
Section & Marks awarded & Maximum & Comment \\
\midrule
Question 1: Changing population & 10 & 10 & Fully correct graph-reading and well-developed explanations of megacity benefits and fertility decline. \\
Question 2: Global climate & 10 & 10 & Accurate definition, strong albedo explanations, and two valid globalization links to increased emissions. \\
Question 3: Resource consumption and security & 10 & 10 & Clear distribution, food-security and energy-security answers with the needed cause-effect links. \\
Question 4: Patterns in population change (Australia) & 8 & 10 & The short answers are correct. The six-mark appraisal has a valid judgement and some named evidence, but needs more distinct supported points from across the infographic. \\
Question 5: Global middle class and water essay & 7 & 10 & Structured and relevant, with clear water-demand pathways and some evaluation against energy; more place-specific evidence and wider comparison are needed for the top band. \\
\midrule
\textbf{Total} & \textbf{45} & \textbf{50} & Main targets: make six-mark answers count five separate evidence points, and strengthen Section C with named examples, statistics and more balanced comparison. \\
\bottomrule
\end{tabularx}

\section*{Question-by-Question Marks}

\begin{longtable}{@{}p{0.12\textwidth} p{0.12\textwidth} p{0.31\textwidth} p{0.37\textwidth}@{}}
\toprule
Question & Mark & Awarded for & Feedback \\
\midrule
\endfirsthead
\toprule
Question & Mark & Awarded for & Feedback \\
\midrule
\endhead
1(a)(i) & \markcell{1}{1} & Gives 18 million for New York in 2000. & Correct; this is within the accepted range. \\
1(a)(ii) & \markcell{1}{1} & Identifies Shanghai. & Correct. \\
1(b) Consequence 1 & \markcell{2}{2} & Explains better healthcare access because megacities attract and concentrate healthcare professionals. & Valid individual benefit with development. \\
1(b) Consequence 2 & \markcell{2}{2} & Links larger population/economic growth and tax revenue to improved public services such as transport. & Valid, developed benefit. To make it even sharper, state explicitly how the individual gains, for example lower commuting cost or improved accessibility. \\
1(c) Reason 1 & \markcell{2}{2} & Links economic growth to higher education, sex education and family planning. & Full credit. Add the final step ``therefore couples have fewer children'' for maximum clarity. \\
1(c) Reason 2 & \markcell{2}{2} & Explains that career and salary priorities can delay childbirth and reduce family size. & Valid and clearly linked to fertility decline. \\
\midrule
2(a) & \markcell{2}{2} & States that anthropogenic greenhouse gases in the atmosphere cause radiation to return to Earth's surface. & Correct. ``Reflect'' is less precise than ``absorb and re-emit/trap'', but the meaning is clear enough for full credit. \\
2(b) Way 1 & \markcell{2}{2} & Explains that warming melts high-albedo ice into lower-albedo water, reducing overall albedo. & Valid human-driven albedo change. \\
2(b) Way 2 & \markcell{2}{2} & Explains that urbanization replaces lighter natural surfaces with darker concrete/asphalt, lowering albedo. & Correct and well developed. \\
2(c) Way 1 & \markcell{2}{2} & Links international trade and imported oil use to greater emissions. & Valid globalization-emissions pathway. A sharper version would say that trade lets countries without domestic oil reserves import and burn oil, increasing greenhouse gas emissions. \\
2(c) Way 2 & \markcell{2}{2} & Links greater global connectivity to more flights and fuel combustion. & Correct. \\
\midrule
3(a) & \markcell{2}{2} & Describes food emergency areas in the south and a smaller northern area near Eritrea. & Correct distribution description. \\
3(b) Factor 1 & \markcell{2}{2} & Explains that lack of fertile land makes agriculture difficult and limits food access. & Valid factor and consequence. \\
3(b) Factor 2 & \markcell{2}{2} & Explains poor storage/distribution infrastructure leading to spoilage or uneven food distribution. & Valid and clearly developed. \\
3(c) Environmental issue & \markcell{2}{2} & Explains that drought/water stress can reduce thermoelectric power operation and hydroelectricity generation. & Full cause-effect link. \\
3(c) Geopolitical issue & \markcell{2}{2} & Explains that conflict over oil-rich land can damage wells, refineries and pipelines. & Valid geopolitical disruption to energy security. \\
\midrule
4(a)(i) & \markcell{1}{1} & Gives New South Wales. & Correct. \\
4(a)(ii) & \markcell{1}{1} & Gives 209,700. & Correct; accepted as the overall range. \\
4(b) & \markcell{2}{2} & Suggests the categories may imply the listed reasons add to 100\%, leaving out other possible migration reasons. & Valid inaccuracy with development. \\
4(c) & \markcell{4}{6} & Gives a support point about movement to main cities, names states/territories with net losses and gains, identifies different capital-city growth percentages using Darwin and Hobart, and gives an overall judgement that the pattern is partly similar but varies strongly by place. & This is a sound appraisal but not yet full. The named state evidence is credited, but it mostly develops one idea about migration direction. For 6/6, include five separate evidence points before the final judgement, drawing from different infographic sections such as city growth, migration direction, reasons for migration, and changes over time. \\
\midrule
5 & \markcell{7}{10} & Defines the global middle class, explains higher demand for virtual water, appliances/infrastructure and meat, compares water with energy, and reaches a relevant conclusion about water being most threatened while resources are interconnected. & This reaches the 7--8 band because it is structured, question-focused and evaluative. It stays at 7 because evidence is still general: add named places, precise statistics, and comparison with at least food and energy. The energy comparison also needs more nuance, since renewable potential does not automatically mean energy security or access is solved. \\
\bottomrule
\end{longtable}

\section*{Teaching Feedback}

\textbf{What is working well}
\begin{itemize}
  \item Two-mark answers usually have both parts: the factor or process, then the explanatory link.
  \item Section A knowledge is secure across population, climate, food security and energy security.
  \item The student reads simple data accurately and uses appropriate geographical vocabulary such as albedo, fertility, infrastructure, virtual water and energy security.
  \item The Section C essay has a clear line of argument rather than just listing facts.
\end{itemize}

\textbf{Main improvement targets}
\begin{itemize}
  \item In six-mark infographic appraisals, count distinct evidence points. The target structure is five data-supported points plus one overall judgement.
  \item In Section C, use named places and statistics. A strong essay should not rely only on global generalizations.
  \item When evaluating a statement, compare alternatives directly. For Question 5, water should be weighed against food, energy and other drivers such as population growth and climate change.
  \item Make final judgement wording precise: say whether you agree fully, partly, or only in specific contexts.
\end{itemize}

\section*{Suggested Next Practice}

\begin{itemize}
  \item Rewrite Question 4(c) as six numbered sentences: three showing similarity, two showing difference, and one final overall judgement.
  \item Build a Section C evidence bank with one named example each for water scarcity, rising meat demand, energy demand and resource management.
  \item Redraft the Question 5 conclusion so it directly compares water, food and energy, then explains why one is the greatest threat in a specific context.
  \item Practise adding one precise statistic or place detail to every essay paragraph.
\end{itemize}

\section*{Model Improvements}

\textbf{Question 4(c) stronger route:}
Use separate evidence points from the infographic before judging the whole pattern. For example: one point on movement to main cities, one on all capital cities increasing, one on states/territories gaining or losing interstate migrants, one on variation in city growth percentages, and one on changing reasons for migration. Then finish with a judgement about whether similarities or differences dominate.

\textbf{Question 5 stronger route:}
The essay should keep the current argument about middle-class consumption and virtual water, but add named evidence. A higher-band version could compare water demand from meat and household appliances with food demand and energy demand, then evaluate whether climate change, population growth or affluence is the larger threat in different places.

\section*{Notes}

\begin{itemize}
  \item The marked student-work source is \texttt{Geography\_paper\_2\_question\_booklet\_TZ1\_HLSL-2025-student-work-1.jpg} through \texttt{student-work-10.jpg}.
  \item The student answered Section C Question 5. Question 6 was not attempted and was not marked.
  \item This report records marking judgements and teaching feedback only; it does not reproduce the exam paper or markscheme.
\end{itemize}

\end{document}
